%% kapa: exam=1 round=26 year=2015 period=1 form=A subject=economics subject_ko=경제원론 seq=18 number=58
%% source: Q-Net 감정평가사 기출문제 (cst00309) artlSeq=5138001 file=extracted/1008/5138001/제26회_감정평가사 1차_기출/제26회_감정평가사 1차_기출/1교시_감평_A형.pdf
%% conversion: kordoc 4.24.51 -> tools/gen1.py (auto); review=proofread-sonnet; flags=none
\begin{question}{58}
A국, B국, C국의 소득분위별 소득점유비중이 다음과 같다. 소득분배에 관한 설명으로 옳은 것은? (단, 1분위는 최하위 20 \%, 5분위는 최상위 20 \%의 가구를 의미한다.)
\hfill(단위: \%)\par
{\small\centering\begin{tabular}{|c|c|c|c|}\hline
 & A국 & B국 & C국 \\ \hline
1분위 & 0 & 20 & 6 \\ \hline
2분위 & 0 & 20 & 10 \\ \hline
3분위 & 0 & 20 & 16 \\ \hline
4분위 & 0 & 20 & 20 \\ \hline
5분위 & 100 & 20 & 48 \\ \hline
\end{tabular}\par}
\begin{choices}
  \item A국은 B국보다 소득분배가 상대적으로 평등하다.
  \item B국은 C국보다 소득분배가 상대적으로 불평등하다.
  \item C국의 십분위분배율은 1/8이다.
  \item A국의 지니계수는 0이다.
  \item B국의 지니계수는 A국의 지니계수보다 작다.
\end{choices}
\end{question}
\answer{5}
